\section{Partition functions up to a polygon boundary}
\label{sec:polygon}

The bulk scale map controls perturbations whose determining neighborhoods
lie in the plane.  We now extend it to a fixed polygon, integrating all
the way to its boundary.  The additional estimates concern logarithms of
partition functions.  The transfer bounds of Section~\ref{sec:boundary}
make the boundary contributions contract without any cancellation in the
environmental mean.

We formulate the comparison on a flat surface because the later slit
argument will use several copies of portions of the same lattice.  This
does not change the local pure model or introduce independent disorder
on the copies.

\begin{definition}[Polygons on sheets]\label{def:sheet-polygon}
Let $\mathcal S$ be an oriented flat surface with a map
$\pi:\mathcal S\to\R^2$ whose restriction to each sufficiently small
neighborhood is an orientation-preserving isometry.  A \emph{collared
orthogonal polygon} is a compact polygonal disk $\mathcal P\subset
\mathcal S$ such that $\pi$ is defined on a neighborhood of
$\mathcal P$, and each side of $\mathcal P$ is horizontal or vertical in
these local coordinates.  The polygon has finitely many sides; the map
$\pi$ need not be injective on $\mathcal P$.

Its lattice approximation is obtained by rounding its sides to the
lifted square grid of mesh $\delta$ and filling the resulting polygon.
Boundary-condition changes are placed in straight side interiors and
rounded with the sides.  Distinct lifted copies of a physical lattice
edge are distinct random-cluster edges with the same environmental
sign.  Thus the coupling of a lifted edge $\widetilde e$ is the coupling
of $\pi(\widetilde e)$.
\end{definition}

The ordinary plane case is included.  A collar permits small changes
of sides and marked positions without changing the local charts.  All
lattice conventions below allow a bounded number of boundary layers;
the rounding prescription is fixed for each test.

\begin{proposition}[Fixed polygon comparison]
\label{prop:polygon-comparison}
There exists $\theta_0>0$ with the following property.  Fix
$0<\theta<\theta_0$, a collared orthogonal polygon $\mathcal P$, and
four distinct boundary marks in straight side interiors.  Wire each of
two opposite marked boundary arcs separately and leave the intervening
arcs free.  On the rounded lifted grid assign odds
\[
 v_e=\sqrt2\exp(\theta\xi_e),
\]
where the physical-edge signs are independent and symmetric, and
repeated lifted edges have the same sign.  Let $\phi_{\delta,J}^{\rm sep}$
and $\phi_{\delta,0}^{\rm sep}$ denote respectively this random-cluster
law and the pure law on the same graph and with the same boundary
partition.  If $\mathcal A_\delta$ is the event of an actual open path
between the two wired arcs, then
\begin{equation}\label{eq:polygon-comparison}
 \phi_{\delta,J}^{\rm sep}(\mathcal A_\delta)
 -\phi_{\delta,0}^{\rm sep}(\mathcal A_\delta)
 \longrightarrow0
 \qquad\text{in environmental probability as }\delta\downarrow0.
\end{equation}
The disorder threshold is independent of the fixed polygon and marks.
The conclusion also holds on the dual square grid.
\end{proposition}

For the proof, constrain the spin to be constant on each wired arc.
Write $\mathcal Z_J^{++}$ and $\mathcal Z_J^{+-}$ for the two partition
functions obtained by fixing these constants, with free spins on the
remaining boundary arcs.  The Edwards--Sokal correspondence gives
\begin{equation}\label{eq:polygon-spin-ratio}
 \phi_{\delta,J}^{\rm sep}(\mathcal A_\delta)
 =\frac{\mathcal Z_J^{++}-\mathcal Z_J^{+-}}
        {\mathcal Z_J^{++}+\mathcal Z_J^{+-}}
 =\tanh\!\left(\frac12\log
                 \frac{\mathcal Z_J^{++}}{\mathcal Z_J^{+-}}\right).
\end{equation}
Indeed, after summing the two possible signs of each arc, their spin
correlation is the probability that their FK clusters are actually
connected.  Global spin flip identifies the other two partition
functions with these two.  We shall compare $\mathcal Z_J^\eta$ with its
pure counterpart $\mathcal Z_0^\eta$, for $\eta\in\{++, +-\}$, while
removing exactly the same scalar factor from both.

\subsection{The scale map with the boundary present}

We first describe the finite-volume implementation sufficiently
explicitly to keep track of that common factor.  Throughout, the norm
of a potential is its supremum over all its spin and reference-variable
arguments.  Independent priorities and thresholds used to construct the
map belong to the environment over which $\E$ averages.  The fair
comparison coins described below are instead Gibbs-field variables,
with an independent product reference law.  We write $\E_0^\eta$ for
pure Ising expectation with pin assignment $\eta$, together with all
these reference variables.

Write $C=C_{\rm s}$ for the fixed support and dependence constant of
Proposition~\ref{prop:bulk-input}, as supplied by the local scale map in
\cite[Section~``An exact local change of scale'']{Q}.  The coarse
spacings are $s_k=s_0L^k$ in lattice units, with $L$ a sufficiently large
odd integer.  All support, allocation, and status-check radii are fixed
multiples of the current spacing.  Their constants can be fixed before
$L$ is chosen.  Choose a positive physical scale $\alpha$, depending
on the polygon, so small that every full determining neighborhood
through scales $\delta s\le\alpha$ lies in an ordinary injectively
projected chart.  If it meets the polygon boundary, it sees at most a
straight side, one corner, or one boundary-condition change.  This is
possible because there are finitely many sides and distinct changes,
and because the closed polygon has a collar.  At a corner whose two
incident rays are pinned, both rays carry the same sign.

For sufficiently small $\delta$, we stop at the largest $N=N(\delta)$ for which
$\delta s_N\le\alpha$.  In particular,
\begin{equation}\label{eq:polygon-terminal-scale}
 \alpha/L<\delta s_N\le\alpha,
 \qquad N(\delta)\longrightarrow\infty.
\end{equation}
The choice of $\alpha$ includes all enlarged neighborhoods used below,
not just the supports of the potentials.  A fixed exterior pad contains
the necessary grid cells and reference variables.  There are no spins
to integrate outside the polygon.  The total overhang accumulated up to
scale $s_N$ is at most a fixed multiple of $s_N$, since the spacings
form a geometric progression.  Consequently an outer truncation beyond
this pad has no effect on any nonzero potential.

Cell indices distinguish different lifts even when their projected
centers coincide.  Within any ordinary chart we use the integer
coordinates of the projected grid.  At spacing $s$, let $f_i$ be the current child potentials and let
$\mathcal C_j$ be the $L^2$ children of the parent at spacing $s'=Ls$.
Set
\[
 \mathcal C_j=Lj+\{-(L-1)/2,\ldots,(L-1)/2\}^2,
 \qquad F_j=\sum_{i\in\mathcal C_j}f_i,
\]
\[
 I_j=\{x:\norm{x-s'j}_\infty<4s'\}.
\]
Coordinates in this display are the ordinary chart coordinates.
Integrating $I_j$ means integrating the free polygon spins and old
reference variables lying there.  The support of $F_j$ lies well inside
$I_j$.

Here is the regular operation of~\cite{Q}.  Fix a numerical range
$D_0=100$ and attach a fair reference coin to every pair of distinct
parents at sup-norm distance at most $D_0$ in parent units.  A parent is
selected when every incident coin favors it. The incident comparisons
form its \emph{comparison star}. Let $A_j$ be its selection
indicator and
\[
 p_{\rm sel}=\E_{\rm ref}A_j=2^{-((2D_0+1)^2-1)}.
\]
Thus $p_{\rm sel}>0$ is fixed independently of $L$, and distinct selected squares
are separated.  Order comparisons by length and then by independent
continuous priorities.  For a comparison edge $e$, let $d_eA_j$ be the
Doob difference obtained by revealing its coin under the reference
law.  This is a local function of the star at $j$: it vanishes if an
earlier comparison has already disfavored $j$.

For fixed old field variables, decorate the new coins successively by
normalized likelihoods
\begin{equation}\label{eq:polygon-decoration}
 \frac{e^{U_e}}{\E_e e^{U_e}},
 \qquad U_e=p_{\rm sel}^{-1}\sum_j F_jd_eA_j,
\end{equation}
where $\E_e$ averages the current fair coin, holding previously revealed
coins fixed.  Since each likelihood is normalized, the decoration
preserves the old field marginal. Only finitely many likelihoods are
nontrivial in this padded finite construction; the plane array is used
to specify their local formulas. The identity
$\sum_e d_eA_j=A_j-p_{\rm sel}$ gives
\begin{equation}\label{eq:polygon-decoration-identity}
 \sum_jF_j+\sum_e\bigl(U_e-\log\E_e e^{U_e}\bigr)
 =p_{\rm sel}^{-1}\sum_j A_jF_j-\sum_e\log\E_e e^{U_e}.
\end{equation}
Every comparison star affecting a nonzero term is retained in full,
including in the exterior pad.  In particular, a normalizer whose old-field support meets an
updated square remains in the expression being integrated.

Given the coins, each normalizer on the right of
\eqref{eq:polygon-decoration-identity} meets at most one selected
square. This is the local separation property of
\cite[Lemma~3.1, \emph{Separation of the terms to be integrated}]{Q}.
Its argument is exactly local: if the endpoint
neighborhoods of a comparison could meet two selected squares, then
for a short comparison the selected centers would be too close; for a
long comparison an endpoint has already lost its shorter comparison
with the nearby selected center.  Its Doob difference is then zero.
The same assertion holds for each selected $F_j$.  If $H_j$ is the sum
of all terms meeting a selected $I_j$, replace it by
\begin{equation}\label{eq:polygon-conditional-addition}
 \log\E_0^\eta[e^{H_j}\mid I_j^c].
\end{equation}
The notation $I_j^c$ includes all unintegrated spins and old reference
variables; the new coins are held fixed.  Pure conditional laws in the
selected squares are independent given their exterior, also when a
square intersects the polygon boundary.  Thus these replacements
preserve the integrated weight exactly.  They leave all old variables
in the field, with pure conditional law in each updated square.

Large potentials are handled by the real-valued modification of the
same map.  Attach to each child an independent environmental threshold
uniform on $[b_0,2b_0]$, and call the child a flag if its norm exceeds
that threshold.  A parent within $10D_0$ parent units of a flag is
inactive.  Set its $F_j$ to zero in~\eqref{eq:polygon-decoration}, retain
its child terms separately, and keep its comparison coins in the full
stars.  Every retained term that touches an active selected square is
included in that square's integration.  Its small support prevents it
from touching two such squares.

Put $d_0=4C+10$.  An isolated group is a nonempty group $A$ of flags of
child-grid diameter at most $6d_0$, with no other flag within $100D_0$
parent units of $A$.  Choose a representative by independent
priorities, let $j(A)$ be its parent, and set
\[
 G_A=\sum_{\operatorname{dist}(i,A)\le3d_0}f_i,
 \qquad I_A=I_{j(A)}.
\]
For large $L$, $G_A$ is supported deep inside $I_A$.  The inactive region
keeps all active integrations and nonzero active coin normalizers away
from $I_A$.  Let $B_A$ be the sum of the remaining terms touching $I_A$.
Replace $G_A$, retaining $B_A$, by
\begin{equation}\label{eq:polygon-isolated-ratio}
 R_A=\log\E_0^\eta[e^{B_A+G_A}\mid I_A^c]
       -\log\E_0^\eta[e^{B_A}\mid I_A^c].
\end{equation}
The isolated groups are separated.  The identity
\begin{equation}\label{eq:polygon-isolated-exactness}
 \E_0^\eta[e^{B_A+R_A}\mid I_A^c]
 =\E_0^\eta[e^{B_A+G_A}\mid I_A^c]
\end{equation}
shows why the background $B_A$ must be kept in this exact operation.
Its removal will be used only in estimating the norm.  Allocate each
integrated term to its integration parent, allocate retained terms to
their parents, and split the remaining pair terms between their
endpoints, as in the fixed local rules of~\cite{Q}.

We next specify the scalar subtractions.  Fix $B$ larger than all the
support and determining-range constants, in parent units, and call a
cell at spacing $s$ a boundary cell if its center is within $Bs$ of the
polygon boundary.  The constant $B$ is fixed before $L$.  At every
scale, including the initial one, use the full-plane pure centering
from~\cite{Q} outside the boundary cells.  For a boundary-cell potential
in the $++$ construction, subtract its $\E_0^{++}$ expectation.  Subtract
that same scalar from the corresponding potential in the $+-$
construction.

\begin{lemma}[Exact boundary implementation]
\label{lem:boundary-exactness}
For the terminal scale chosen above, the local construction is exact
for each fixed environment and realization of its priorities and
thresholds.  Its potentials have support and environmental dependence
radius at most $Cs_k$.  Outside the
boundary cells they are the plane potentials, or zero outside the pad.
The two pin assignments admit the same scalar subtractions.  More
precisely, there is a scalar $a_\delta$, independent of the pin
assignment, such that
\begin{equation}\label{eq:polygon-exact-partition}
 \frac{\mathcal Z_J^\eta}{\mathcal Z_0^\eta}
 =e^{a_\delta}\E_0^\eta
       \exp\!\left(\sum_i f_i^{(N),\eta}\right),
 \qquad \eta\in\{++, +-\}.
\end{equation}
The initial boundary norms are uniformly $O_{s_0}(\theta)$.
\end{lemma}

\begin{proof}
Allocate the initial bond perturbations as in
\cite[Section~``Tuning a trajectory with vanishing perturbations'']{Q},
using only bonds present in the polygon.  A cell receives the allocated
terms $(K_e-K_0)\sigma_x\sigma_y$, followed by the specified centering.
For fixed $s_0$ there are boundedly many such terms per cell, and
$|K_e-K_0|=O(\theta)$; this proves the initial bound.  Subtracting
constants changes only the scalar partition factor.

The identities \eqref{eq:polygon-decoration-identity},
\eqref{eq:polygon-conditional-addition}, and
\eqref{eq:polygon-isolated-exactness} prove exactness at each step.
They hold for every conditional pure probability kernel; the presence
of the boundary merely reduces the set of spins integrated in a
square.  The allocation and status checks use the same finite stencils
as in the plane.  The choice of $C$ in the local map already covers
these stencils and the $C/L$ overhang from the child scale.  Restricting
the pure integrations to polygon spins does not enlarge either radius.
This proves locality.  Choosing
$B$ with slack so that $B>C_1+B/L$, where $C_1$ bounds the complete
parent determining radius, shows inductively that every input of a
nonboundary output is also nonboundary.  There the construction is
precisely the plane construction.  The pad similarly contains all
possibly nonzero outputs.

Every determining neighborhood sees at most one local pin sign.
Spin flip commutes with the conditional integrals, and with every
algebraic operation in the decoration and the isolated ratio.  Hence
changing that sign conjugates the local potential by simultaneous
flip of its spin and pin arguments.  The corresponding norms agree;
using the same thresholds and priorities therefore gives the same
flags, isolation decisions, and local scalar subtraction.  If no pin
lies in a determining neighborhood, its potential is spin-flip even.
The constants used for boundary centering are deterministic
functionals of the local potential: dependence on the fixed pure
polygon does not introduce dependence on remote random bonds.

In particular the $+-$ potentials need not be centered for the pure
$+-$ expectation.  Exactness only requires that the subtracted
constants agree between the two constructions.  Iterating the exact
identities with those constants gives
\eqref{eq:polygon-exact-partition}.
\end{proof}

\subsection{Contraction of the boundary potentials}

Let $X_k$ be the root moment of the plane potentials supplied by
Proposition~\ref{prop:bulk-input}, with its fixed order $P>2$, and let
\begin{equation}\label{eq:polygon-boundary-moment}
 Z_k=\max_{i\text{ a boundary cell}}
       \bigl(\E\norm{f_i^{(k),++}}_\infty^P\bigr)^{1/P}.
\end{equation}
The maximum is zero if there are no boundary cells.  By local spin-flip
conjugacy the same quantity bounds the $+-$ construction.  We take the
maximum of the individual moments, rather than the moment of a maximum.
No stationarity along the polygon boundary is required.

\begin{lemma}[Boundary recurrence]\label{lem:boundary-recursion}
The scale-step parameters can be chosen in the order: $L$ sufficiently
large, $b_0$ sufficiently small, and $s_0$ sufficiently large, so that
for all $k<N$,
\begin{equation}\label{eq:boundary-recurrence}
 Z_{k+1}\le q_* Z_k+C_LX_k
                    +C_{L,b_0}(Z_k+X_k)^2,
 \qquad q_*<1.
\end{equation}
The constants are independent of the polygon and the mesh, provided
the determining neighborhoods have the boundary geometries and chart
properties specified above.  For sufficiently small $\theta$, the
boundary moments remain uniformly small and tend to zero as the
scale index tends to infinity, uniformly over all the finite runs.
\end{lemma}

\begin{proof}
There are four contributions to the recurrence.  The linear transfer
of boundary logs contracts by the summable side estimate; unflagged
Taylor remainders are small because their input norms are at most
$2b_0$.  On flagged stencils, a distant independent flag gives an
extra moment factor.  The remaining dependent logs form an isolated
group, to which the positive-weight transfer estimate applies.
Bulk inputs in each case are bounded by a constant times $X_k$.

We first estimate the linearized regular map.  Its plane formula from
\cite[Lemma~3.2, \emph{Regular linearization}]{Q} remains valid before the
specified output centering:
\[
 p_{\rm sel}^{-1}A_j\sum_{i\in\mathcal C_j}
          \E_0^{++}[f_i\mid I_j^c].
\]
The derivative of each coin normalizer is zero, and the derivative of
a conditional log integral is conditional expectation.  For a
boundary child the pure mean is zero, so its transferred supremum is
bounded by its oscillation.  The reference variables in its patch are
integrated with their pure product law.  The same bound for the other
pin assignment follows by the local flip described in
Lemma~\ref{lem:boundary-exactness}.

We claim that the sum of the transfer coefficients over boundary
children in one determining stencil is $o_L(1)$.  Distances in the next
paragraph are measured in child units.  All these cells lie in a strip
of fixed width $B$ along a side.  If a cell is within a fixed distance
of a corner or a change, choose an inner sector shell containing its
whole support and determining patch.  The corner or change estimate
of Lemma~\ref{lem:boundary-transfer}, composed to a radius proportional
to $L$, makes its transfer coefficient tend to zero as $L$ grows.
All those shells fit inside the parent integration square with room
to condition onward.

For the remaining cells let $t$ be their distance along the side to
the nearest corner or change, capped at $L$; in a neighborhood without
a feature set $t=L$.  Once $t$ exceeds a fixed support-dependent
constant, a straight half-shell reaches a radius proportional to $t$
without seeing a different boundary condition.  Its coefficient is
at most $C't^{-2}$ by Lemma~\ref{lem:boundary-transfer}.  On a free side
the full determining stencil avoids pins, so the potential is even,
as required by that lemma.  There are $O(t)$ boundary cells in each
dyadic range $[t,2t]$, including $O(L)$ in the last capped range.
Consequently, for any sufficiently large fixed $t_0$,
\begin{equation}\label{eq:boundary-summable-transfer}
 \sum_i a_{i,L}
 \le \sum_{t_i\le t_0}a_{i,L}+\frac{C''}{t_0},
 \qquad
 \lim_{L\to\infty}a_{i,L}=0
       \quad\text{for each member of the first, bounded list}.
\end{equation}
First choosing $t_0$ and then $L$ proves the claim.  The same calculation
holds for a fixed enlargement of the child patches or the boundary
strip.  Patches with no free spins cause no difficulty, since only
their reference variables remain to be integrated.

Multiplication by $A_j/p_{\rm sel}$ costs at most $p_{\rm sel}^{-1}$, which is independent
of $L$, and output centering costs at most a factor two.  Minkowski's
inequality and \eqref{eq:boundary-summable-transfer} therefore bound
the boundary-child part of the linearized output by $o_L(1)Z_k$.
All bulk inputs can be bounded crudely, giving $C_LX_k$.
Only finitely many fixed shell ratios are needed once $L$ is fixed;
we subsequently choose $s_0$ large enough for all their pure estimates.

It remains to control the actual map in the $P$th moment.  Set
$X=Z_k+X_k$.  Every child norm in an enlarged determining stencil has
root moment at most $X$.  On a stencil without flags, every child norm
is at most $2b_0$.  Conditional integrations are probability kernels,
so the regular map has uniform Taylor bounds after $L$ is fixed.
Writing $S$ for the sum of its child norms, its remainder is bounded
by $C_LS^2\le C_Lb_0S$.  Its root-moment cost is thus $C_Lb_0X$.
The real map, without any smallness restriction, has the crude bound
$C_LS$ by $|\log\E_0e^H|\le\norm{H}_\infty$ and the analogous
Lipschitz bound for differences of such logarithms.

Small flag probability alone would not improve this last estimate in
the $P$th moment.  We use the location of the flags, as in
\cite[Section~``The contribution of isolated large terms'']{Q}.
Enlarge the determining stencil to include all isolation checks and
all child neighborhoods of radius $100D_0L+2C$ about a contributing
child.  If a flag $h$ in this enlarged stencil is farther than $2C$
child units from a charged child $i$, their environmental dependence
neighborhoods are disjoint.  Its fresh threshold is independent too,
and Markov's inequality gives
$\P(h\text{ is a flag})\le b_0^{-P}X^P$.  Hence
\begin{equation}\label{eq:boundary-independent-flag}
 \E\!\left[\norm{f_i}_\infty^P
                    {\bf1}_{\{h\text{ is a flag}\}}\right]
 =\E\norm{f_i}_\infty^P\,
       \P(h\text{ is a flag})
 \le C_{b_0}X^{2P}.
\end{equation}
There are only finitely many charges for fixed $L$, so these
contributions cost $C_{L,b_0}X^2$ in root moment.

If no such independent flag is available for a contributing term,
every relevant flag lies within $2C$ of that term.  These flags form
one bounded child-radius group.  The enlarged stencil contains the
whole isolation neighborhood, so either they are an isolated group
of the construction, or a farther flag is present and
\eqref{eq:boundary-independent-flag} applies.  In the isolated case,
all logs possibly dependent on the triggering flags are included in
$G_A$.  For every probability kernel $P_0$,
\begin{equation}\label{eq:boundary-remove-background}
 \left|\log\frac{P_0e^{B_A+G_A}}{P_0e^{B_A}}
                    -\log P_0e^{G_A}\right|
 \le2\norm{B_A}_\infty.
\end{equation}
Each child in $B_A$ is outside the concentration neighborhood and is
independent of an actual triggering flag.  Its occurrence can again be
charged using \eqref{eq:boundary-independent-flag}.  The same applies
to every term omitted from the concentrated contribution.

After this estimate, the only remaining large term to bound is the
centered form of
\[
 \log\E_0^{++}[e^{G_A}\mid I_A^c].
\]
Corollary~\ref{cor:boundary-observables} applies to this positive
weight for any size of $\norm{G_A}_\infty$.  A group containing a
boundary child has bounded child diameter.  There are boundedly many
possible flag patterns at each anchor and boundedly many logs in each
concentration neighborhood, with bounds independent of $L$.
Thus its logs all lie in
a fixed enlargement of the boundary strip, and the shell coefficient
at its anchor multiplies the sum of their norms.  Summing over all
possible boundary anchors gives the same arbitrarily small sum as in
\eqref{eq:boundary-summable-transfer}.  We may sum before taking the
root moment by Minkowski's inequality; neither a supremum over anchors
nor independence among their events is needed.  Groups containing
only bulk inputs contribute at most $C_LX_k$ by the crude bound.
The linearized terms restricted to flagged stencils, if subtracted in
comparing the two maps, obey the identical dichotomy: either charge
them to an independent flag or use the linear shell transfer for the
bounded local group.

Combining these estimates gives
\[
 Z_{k+1}\le
  \bigl(o_L(1)+C_Lb_0\bigr)Z_k+C_L'X_k
                   +C_{L,b_0}(Z_k+X_k)^2.
\]
Choose $L$ so the first term leaves a strict contraction margin, then
choose $b_0$ within that margin, and finally choose $s_0$ for the
finite list of pure shell estimates.  These choices are compatible
with the parameter freedoms in Proposition~\ref{prop:bulk-input}.
This proves \eqref{eq:boundary-recurrence}.

For completeness, take a deterministic scalar majorant for all the
finite polygon runs, starting from their uniform $O_{s_0}(\theta)$
bound and iterating the right side of
\eqref{eq:boundary-recurrence}.  Choose a small invariant interval so
that the quadratic term has Lipschitz constant less than half the
remaining contraction margin.  Proposition~\ref{prop:bulk-input}
makes $\sup_kX_k$ as small as required by decreasing $\theta$; the
initial majorant is small by the same choice.  The majorant therefore
stays in that interval.  Since $X_k\to0$, its limsup $z$ satisfies
$z\le q_*z+C_{L,b_0}z^2$, which forces $z=0$ in the chosen interval.
This proves uniform smallness and decay.
\end{proof}

\subsection{Completion of the partition comparison}

\begin{proof}[Proof of Proposition~\ref{prop:polygon-comparison}]
All determining stencils used in Lemma~\ref{lem:boundary-recursion}
lie in ordinary injectively projected charts.  The physical-edge
signs within each such stencil have their usual iid law.  In
particular, disjoint dependence neighborhoods used in
\eqref{eq:boundary-independent-flag} are genuinely independent.
There is no requirement of independence between different sheets:
the entire enlarged stencil, including its overhang, was placed in
one ordinary chart before the terminal scale was selected.

By \eqref{eq:polygon-terminal-scale}, the number of possibly nonzero
terminal potentials, including those in the pad, is bounded for this
fixed polygon.  This bound may depend on its geometry and on
$\alpha$, but not on the mesh.  Proposition~\ref{prop:bulk-input} and
Lemma~\ref{lem:boundary-recursion} therefore give
\[
 S_\delta=\sum_i\max_{\eta\in\{++, +-\}}
                 \norm{f_i^{(N),\eta}}_\infty
 \longrightarrow0
 \quad\text{in }L^1\text{ of the environment and construction choices}.
\]
Applying \eqref{eq:polygon-exact-partition} yields
\[
 \left|\log\frac{\mathcal Z_J^\eta}{\mathcal Z_0^\eta}
                      -a_\delta\right|\le S_\delta,
 \qquad \eta\in\{++, +-\}.
\]
The common scalar cancels from the difference.  Since $\tanh$ is
Lipschitz, \eqref{eq:polygon-spin-ratio} proves the comparison with an
error bounded by $S_\delta$.  The partition ratios themselves do not
depend on the construction choices.  Averaging the bound over those
choices thus gives convergence in environmental probability alone.

The local boundary shapes are the straight side, the two orthogonal
corner angles, and a single fixed/free change.  Their constants, and
all scale-map parameters, were chosen independently of the polygon;
only the positive terminal physical scale depended on it.  This
proves the asserted uniform disorder threshold.  On the dual grid,
duality replaces $v_e$ by $2/v_e$, hence $\theta\xi_e$ by
$-\theta\xi_e$.  This has the same environmental law, and the same
argument applies.
\end{proof}
